\documentclass{article}
\usepackage{graphicx}
\usepackage{amsmath}
\usepackage{amsfonts}
\usepackage{bookmark}
\usepackage{booktabs}
\usepackage{exercises}
\graphicspath{ {} }
\usepackage{indentfirst}
\usepackage{amsthm}
\usepackage{yhmath}
\usepackage{geometry}
\usepackage{placeins}
\usepackage{csquotes}
\usepackage{pgfplots}
\newcommand*{\oldneg}{\mathord{\sim}}
\newcommand\showdiv[1]{\overline{\smash{)}#1}}
\newcommand\Mydiv[1]{%
\smash{\raise.3ex\hbox{$\big)$}}\mkern-8mu
        \overline{\enspace\strut#1}}
\setlength{\tabcolsep}{5pt}
\theoremstyle{definition}
\newtheorem{definition}{Definition}

\setlength{\parindent}{.5in}
\setlength{\parskip}{12pt}
\renewcommand{\arraystretch}{2}
\geometry{margin=1in}
\newcommand{\?}{\stackrel{?}{=}}
\pgfplotsset{width=7cm,compat=1.18}\usepgfplotslibrary{statistics}

\title{10.7 Answers}
\begin{document}
\maketitle

\section*{Derivations of Half and Double Angle Formulas}
The double angle formulas follow from the sum formula:
\[ \sin 2a=\sin (a+a)=\sin a\cos a+\sin a\cos a=2\sin a\cos a \]
and
\[ \tan 2a=\tan (a+a)=\frac{\tan a+\tan a}{1-\tan a\tan a}=\frac{2\tan a}{1-\tan^2 a} \]
and
\[ \cos 2a=\cos (a+a)=\cos a\cos a-\sin a\sin a=\cos^2 a-\sin^2 a \]
Now, by the Pythagorean Identity, $\sin^2 a=1-\sin^2 a$, and $\cos^2 a=1-\sin^2 a$. Thus
\[ \cos^2 a-\sin^2 a=1-\sin^2 a-\sin^2 a=1-2\sin^2 a \] and
\[ \cos^2 a-\sin^2 a=\cos^2 a-(1-\cos^2 a)=2\cos^2 a-1 \]
To derive the half angle formulas, we need only to derive the one for cosine (then the rest will follow). This can be done by considering the diagram in Problem 49. The derivation can be 
done by using the Law of Cosines. Consider the unknown hypotenuse $h$ of the right triangle. We can find it using the Pythagorean Theorem:
\[ h^2=(1+\cos\theta)^2+\sin^2\theta \]
\[ =1+\cos^2\theta+2\cos\theta+\sin^2\theta \]
\[ =2+2\cos\theta\textrm{ (By the Pythagorean Identity)} \]
meaning
\[ h=\sqrt{2+2\cos\theta} \]
Note that, because $\theta$ is acute, $h=+\sqrt{2+2\cos\theta}$. Then we can use the law of cosines to find $\cos\frac{\theta}{2}$:
\[ \cos\frac{\theta}{2}=\frac{1^2+h^2-1^2}{2(1)(h)} \]
\[ =\frac{h^2}{2h} \]
\[ =\dfrac{2+2\cos\theta}{2\sqrt{2+2\cos\theta}} \]
\[ =\dfrac{1+\cos\theta}{\sqrt{2+2\cos\theta}} \]
\[ =\dfrac{1+\cos\theta}{\sqrt{2(1+\cos\theta)}} \]
\[ =\dfrac{1+\cos\theta}{\sqrt{2}\sqrt{1+\cos\theta}} \]
\[ =\dfrac{\sqrt{1+\cos\theta}}{\sqrt{2}}\textrm{ (Because $\frac{1+\cos\theta}{\sqrt{1+\cos\theta}}=\sqrt{1+\cos\theta}$)} \]
\[ =\sqrt{\frac{1+\cos\theta}{2}} \]
thus completing the proof in this case.


We can also derive the half angle formula for the law of cosines from the double angle formula:
\[ \cos(2\frac{\theta}{2})=2\cos^2\frac{\theta}{2}-1 \]
\[ \cos^2\frac{\theta}{2}=\frac{\cos\theta+1}{2} \]
\[ \cos\frac{\theta}{2}=\pm\sqrt{\frac{\cos\theta+1}{2}} \]
Now we derive the formula for sine
\[ \cos 2\frac{\theta}{2}=\cos^2\frac{\theta}{2}-\sin^2\frac{\theta}{2} \]
\[ \sin^2\frac{\theta}{2}=\cos^2\frac{\theta}{2}-\cos\theta \]
\[ \sin\frac{\theta}{2}=\pm\sqrt{\cos^2\frac{\theta}{2}-\cos\theta} \]
\[ =\pm\sqrt{\frac{1+\cos\theta}{2}-\cos\theta} \]
\[ =\pm\sqrt{\frac{1-\cos\theta}{2}} \]
Lastly, we derive the formula for tangent:
\[ \tan\frac{\theta}{2}=\dfrac{\sin\frac{\theta}{2}}{\cos\frac{\theta}{2}} \]
\[ =\dfrac{\pm\sqrt{\frac{1-\cos\theta}{2}}}{\pm\sqrt{\frac{1+\cos\theta}{2}}} \]
\[ =\pm\sqrt{\dfrac{\frac{1-\cos\theta}{2}}{\frac{1+\cos\theta}{2}}} \]
\[ =\pm\sqrt{\frac{1-\cos\theta}{1+\cos\theta}} \]
Now, from this, we can derive the two formulas for tangent:
\[ \pm\sqrt{\frac{1-\cos\theta}{1+\cos\theta}} \]
\[ =\pm\sqrt{\frac{(1-\cos\theta)(1+\cos\theta)}{(1+\cos\theta)^2}} \]
\[ =\pm\sqrt{\frac{1-\cos^2\theta}{(1+\cos\theta)^2}} \]
\[ =\pm\sqrt{\frac{\sin^2\theta}{(1+\cos\theta)^2}} \]
\[ =\pm\frac{\sin\theta}{1+\cos\theta} \]
and also
\[ \pm\sqrt{\frac{1-\cos\theta}{1+\cos\theta}} \]
\[ =\pm\sqrt{\frac{(1-\cos\theta)^2}{(1+\cos\theta)(1-\cos\theta)}} \]
\[ =\pm\sqrt{\frac{(1-\cos\theta)^2}{\sin^2\theta}} \]
\[ =\pm\frac{1-\cos\theta}{\sin\theta} \]
Note that, in both cases, we always take the positive fraction, not the negative. For suppose $\tan\frac{\theta}{2}$ is positive. Then either $\frac{\theta}{2}$ is in Quadrant I or 
Quandrant III. In the first case, $\theta$ is in Quandrant II. Hence $\cos\theta$ is negative and $\sin\theta$ is positive. Thus both
\[ \frac{\sin\theta}{1+\cos\theta}\textrm{ and }\frac{1-\cos\theta}{\sin\theta} \] are positive, meaning we take the positive fraction. If $\frac{\theta}{2}$ is in Quadrant III, then 
$\theta$ is either in Quadrant I or II. In any case, $\sin\theta$ is positive, and thus we take the positive sign.


Now suppose $\tan\frac{\theta}{2}$ is negative. Thus either $\frac{\theta}{2}$ is in Quadrant II or Quadrant IV. In the first case, either $\theta$ is in Quadrant III or Quadrant IV, meaning 
$\sin\theta$ is negative, and thus so is
\[ \frac{\sin\theta}{1+\cos\theta}\textrm{ and }\frac{1-\cos\theta}{\sin\theta} \] Now, in the second case, when $\frac{\theta}{2}$ is in Quadrant IV, $\theta$ is again in either Quadrant III or 
Quadrant IV, and the same reasoning follows. 

\section*{6}
For this, note that $-150^\circ$ is coterminal with $-150+360=210^\circ$, which has a reference angle of $210-180=30^\circ$ and lies in Quadrant II. Also note that 
$-75^\circ$ lies in Quadrant IV.
\[ \cos(-75^\circ)=\cos(-\frac{150}{2}^\circ) \]
\[ =\pm\sqrt{\frac{1+\cos(-150^\circ)}{2}} \]
\[ =\pm\sqrt{\frac{1-\frac{\sqrt{3}}{2}}{2}} \]
\[ =\pm\sqrt{\frac{\frac{2-\sqrt{3}}{2}}{2}} \]
\[ =\frac{\sqrt{2-\sqrt{3}}}{2} \]
We take the positive square root because $-75^\circ$ lies in Quadrant IV.

\section*{8}
For this, note that $\frac{5\pi}{6}$ is has a reference angle of $\frac{\pi}{6}$ and lies in Quadrant II. Thus
\[ \sin\frac{5\pi}{12}=\sin\dfrac{\frac{5\pi}{6}}{2} \]
\[ =\pm\sqrt{\frac{1-\cos\frac{5\pi}{6}}{2}} \]
\[ =\pm\sqrt{\frac{1+\frac{\sqrt{3}}{2}}{2}} \]
\[ =\sqrt{\frac{1+\frac{\sqrt{3}}{2}}{2}}\textrm{ (Because $\frac{5\pi}{12}$ lies in Quadrant I)} \]
\[ =\frac{\sqrt{2+\sqrt{3}}}{2} \]

\section*{9}
For this, note that $-\frac{5\pi}{4}$ is coterminal with $-\frac{5\pi}{4}+2\pi=\frac{3\pi}{4}$, which lies in Quadrant II and has a reference angle of $\frac{\pi}{4}$. Also note that 
$-\frac{5\pi}{8}$ lies in Quadrant I. Thus
\[ \tan(-\frac{5\pi}{8})=\tan(-\frac{\frac{5\pi}{4}}{2}) \]
\[ =\frac{\sin(-\frac{5\pi}{4})}{1+\cos(-\frac{5\pi}{4})} \]
\[ =\dfrac{\frac{1}{\sqrt{2}}}{1-\frac{1}{\sqrt{2}}} \]
\[ =\dfrac{\frac{1}{\sqrt{2}}}{\frac{\sqrt{2}-1}{\sqrt{2}}} \]
\[ =\frac{1}{\sqrt{2}-1} \]
\[ =\frac{-1-\sqrt{2}}{(-1+\sqrt{2})(-1-\sqrt{2})} \]
\[ =\frac{-1-\sqrt{2}}{-1} \]
\[ =\sqrt{2}+1 \]

\section*{11}
\[ \tan 15^\circ=\tan\frac{30^\circ}{2} \]
\[ =\frac{1+\cos 30^\circ}{\sin 30^\circ} \]
\[ =\frac{1+\frac{\sqrt{3}}{2}}{\frac{1}{2}} \]
\[ =2+\sqrt{3} \]
Thus the answer is D.

\section*{15}
We are given that $\sin a=-\frac{3}{5}$ and $\pi <a<\frac{3\pi}{2}$. By the Pythagorean Identity, we know that:
\[ \sin^2 a+\cos^2 a=1 \]
\[ (-\frac{3}{5})^2+\cos^2 a=1 \]
\[ \frac{9}{25}+\cos^2 a=1 \]
\[ \cos^2 a=\frac{16}{25} \]
\[ \cos a=\pm\frac{4}{5} \]
\[ \cos a=-\frac{4}{5}\textrm{ (Because $a$ lies in Quadrant III)} \]
meaning
\[ \sin\frac{a}{2}=\pm\sqrt{\frac{1-\cos a}{2}} \]
\[ =\pm\sqrt{\frac{1+\frac{4}{5}}{2}} \]
\[ =\pm\sqrt{\frac{9}{10}} \]
\[ =\pm\frac{3}{\sqrt{10}}=\pm\frac{3\sqrt{10}}{10} \]
\[ =\frac{3\sqrt{10}}{10}\textrm{ (Because $\frac{a}{2}$ lies in Quadrant II)} \]
and
\[ \cos\frac{a}{2}=\pm\sqrt{\frac{1+\cos a}{2}} \]
\[ =\pm\sqrt{\frac{1-\frac{4}{5}}{2}} \]
\[ =\pm\sqrt{\frac{1}{10}} \]
\[ =-\frac{\sqrt{10}}{10}\textrm{ (Because $\frac{a}{2}$ lies in Quadrant II)} \]
and
\[ \tan\frac{a}{2}=\frac{1-\cos a}{\sin a} \]
\[ =\frac{1+\frac{4}{5}}{-\frac{3}{5}} \]
\[ =\dfrac{\frac{9}{5}}{-\frac{3}{5}} \]
\[ =-3 \] 

\section*{20}
We are given that $\cos a=\frac{2}{5}$ and $-\frac{\pi}{2}<a<0$. Thus $a$ is in Quadrant IV. By the Pythagorean Identity, we know that:
\[ \sin^2 a+\cos^2 a=1 \]
\[ \sin^2 a+\frac{4}{25}=1 \]
\[ \sin^2 a=\frac{21}{25} \]
\[ \sin a=\pm\frac{\sqrt{21}}{5} \]
\[ \sin a=-\frac{\sqrt{21}}{5}\textrm{ (Because $a$ is in Quadrant IV)} \]
meaning
\[ \tan a=\dfrac{-\frac{\sqrt{21}}{5}}{\frac{2}{5}} \]
\[ =-\frac{\sqrt{21}}{2} \]
and thus
\[ \sin 2a=2\sin a\cos a \]
\[ =2(-\frac{\sqrt{21}}{5})(\frac{2}{5}) \]
\[ =-\frac{4\sqrt{21}}{25} \]
and
\[ \cos 2a=\cos^2 a-\sin^2 a \]
\[ =(\frac{2}{5})^2-(-\frac{\sqrt{21}}{5})^2 \]
\[ =\frac{4}{25}-\frac{21}{25} \]
\[ =-\frac{17}{25} \]
and
\[ \tan 2a=\frac{2\tan a}{1-\tan^2 a} \]
\[ =\dfrac{2(-\frac{4\sqrt{11}}{25})}{1-(-\frac{4\sqrt{21}}{25})^2} \]
\[ =\dfrac{-\frac{8\sqrt{11}}{25}}{1-\frac{336}{625}} \]
\[ =\dfrac{-\frac{8\sqrt{11}}{25}}{\frac{289}{336}} \]
\[ =-\frac{2688\sqrt{11}}{7225} \]

\section*{21}
\[ \frac{\cos 2\theta}{1-2\sin^2\theta} \]
\[ =\frac{\cos 2\theta}{\cos 2\theta}\textrm{ (Double angle formula)} \]
\[ =1 \]

\section*{22}
\[ \frac{\sin 2\theta}{2\cos\theta} \]
\[ =\frac{2\sin\theta\cos\theta}{2\cos\theta} \]
\[ =\sin\theta \]

\section*{23}
\[ (1-\tan\theta)\tan 2\theta \]
\[ =\tan 2\theta-\tan\tan 2\theta \]
\[ =\frac{2\tan\theta}{1-\tan^2\theta}-\frac{2\tan^2\theta}{1-\tan^2\theta} \]
\[ =\frac{2\tan\theta(1-\tan\theta)}{1-\tan^2\theta} \]
\[ =\frac{2\tan\theta(1-\tan\theta)}{(1+\tan\theta)(1-\tan\theta)} \]
\[ =\frac{2\tan\theta}{1+\tan\theta} \]

\section*{24}
\[ \frac{\cos 2\theta}{\sin\theta-\cos\theta} \]
\[ =\frac{\cos^2\theta-\sin^2\theta}{\sin\theta-\cos\theta} \]
\[ =\frac{(\cos\theta+\sin\theta)(\cos\theta-\sin\theta)}{\sin\theta-\cos\theta} \]
\[ =\frac{-(\sin\theta-\cos\theta)(\cos\theta+\sin\theta)}{\sin\theta-\cos\theta} \]
\[ =-\cos\theta-\sin\theta \]

\section*{25}
\[ \dfrac{-\tan\frac{\theta}{2}}{\csc\theta} \]
\[ =\dfrac{-\frac{1-\cos\theta}{\sin\theta}}{\csc\theta} \]
\[ =\cos\theta-1 \]

\section*{26}
\[ 2\sin\frac{\theta}{2}\cos\frac{\theta}{2} \]
\[ \sin\theta\textrm{ (By the double angle formula)} \]

\section*{27}
\[ \cot\theta+\tan\theta \]
\[ =\frac{1}{\tan\theta}+\tan\theta \]
\[ =\frac{1+\tan^2\theta}{\tan\theta} \]
\[ =\frac{\sec^2\theta}{\tan\theta}\textrm{ (Pythagorean Identity)} \]
\[ =\frac{1}{\cos^2\theta\tan\theta} \]
\[ =\frac{1}{\sin\theta\cos\theta}\textrm{ (Tangent Identity)} \]
\[ =2\frac{1}{2\sin\theta\cos\theta} \]
\[ =2\frac{1}{\sin 2\theta} \]
\[ =2\csc 2\theta \]
Thus B is the correct option.

\section*{30}
\[ 2\cos^2\theta \]
\[ =1+(2\cos^2\theta-1) \]
\[ =1+\cos 2\theta \]

\section*{31}
\[ \sin 3\theta=\sin(\theta+2\theta) \]
\[ =\sin\theta\cos 2\theta+\sin 2\theta\cos\theta \]
\[ =\sin\theta(2\cos^2\theta-1)+2\sin\theta\cos^2\theta \]
\[ =2\sin\theta\cos^2\theta-\sin\theta+2\sin\theta\cos^2\theta \]
\[ =4\sin\theta\cos^2\theta-\sin\theta \]
\[ =\sin\theta(4\cos^2\theta-1) \]

\section*{32}
\[ \frac{1}{2}\sin\frac{2x}{3}=\frac{1}{2}2\sin\frac{x}{3}\cos\frac{x}{3} \]
\[ =\sin\frac{x}{3}\cos\frac{x}{3} \]

\section*{33}
\[ 2\sin^2 x\tan\frac{x}{2}=2\sin^2 x\frac{1-\cos x}{\sin x} \]
\[ =2\sin x(1-\cos x) \]
\[ =2\sin x-2\sin x\cos x \]
\[ =2\sin x-\sin 2x\textrm{ (Double Angle Formula)} \]

\section*{34}
\[ -\frac{\cos 2\theta}{\sin\theta}=-\frac{-2\sin^2\theta+1}{\sin\theta} \]
\[ =\frac{2\sin^2\theta-1}{\sin\theta} \]
\[ =\frac{2\sin^2\theta}{\sin\theta}-\frac{1}{\sin\theta} \]
\[ =2\sin\theta-\csc\theta \]

\section*{35}
\[ \cos 4\theta=\cos^2 2\theta-\sin^2 2\theta \]
\[ =(\cos^2\theta-\sin^2\theta)^2-(2\sin\theta\cos\theta)^2 \]
\[ =\cos^4\theta+\sin^4\theta-2\sin^2\theta\cos^2\theta-4\sin^2\theta\cos^2\theta \]
\[ =\cos^4\theta-6\sin^2\theta\cos^2\theta+\sin^4\theta \]

\section*{41}
\[ 2\sin 2x\sin x=3\cos x \]
\[ 4\sin^2 x\cos x=3\cos x \]
\[ \sin x=\pm\frac{\sqrt{3}}{2} \]
In the range $0\le x\le 2\pi$, the solutions are 
\[ \frac{\pi}{3},\pi-\frac{\pi}{3}=\frac{2\pi}{3},\pi+\frac{\pi}{3}=\frac{4\pi}{3},2\pi-\frac{\pi}{3}=\frac{5\pi}{3} \]
We check these solutions:
\[ 2\sin (2\frac{\pi}{3})\sin\frac{\pi}{3}\? 3\cos\frac{\pi}{3} \]
\[ 2\frac{\sqrt{3}}{2}\frac{\sqrt{3}}{2}\? \frac{3}{2} \]
\[ \frac{3}{2}=\frac{3}{2} \]
and
\[ 2\sin (2\frac{2\pi}{3})\sin\frac{2\pi}{3}\? 3\cos\frac{2\pi}{3} \]
\[ 2(-\frac{\sqrt{3}}{2})\frac{\sqrt{3}}{2}\? -\frac{3}{2} \]
\[ -\frac{3}{2}=-\frac{3}{2} \]
and
\[ 2\sin (2\frac{4\pi}{3})\sin\frac{4\pi}{3}\? 3\cos\frac{4\pi}{3} \]
\[ 2\frac{\sqrt{3}}{2}(\frac{-\sqrt{3}}{2})\? 3(-\frac{1}{2}) \]
\[ -\frac{3}{2}=-\frac{3}{2} \]
and
\[ 2\sin (2\frac{5\pi}{3})\sin\frac{5\pi}{3}\? 3\cos\frac{5\pi}{3} \]
\[ 2(-\frac{\sqrt{3}}{2})(-\frac{\sqrt{3}}{2})\? \frac{3}{2} \]
\[ \frac{3}{2}=\frac{3}{2} \]
thus meaning none of the solutions are extraneous.

\section*{43}
\[ \tan\frac{x}{2}=\sin x \]
\[ \frac{1-\cos x}{\sin x}=\sin x \]
\[ 1-\cos x=1 \]
\[ \cos x=0 \]
which occurs when $x=\frac{\pi}{2}$ or $x=\frac{3\pi}{2}$. We check these solutions:
\[ \tan\frac{\frac{\pi}{2}}{2}\?\sin\frac{\pi}{2} \]
\[ 1=1 \]
and
\[ \tan\frac{\frac{3\pi}{2}}{2}\?\sin\frac{3\pi}{2} \]
\[ -1=-1 \]
thus none of the solutions are extraneous, meaning the general solution is
\[ x=\frac{\pi}{2}+\pi n \] where $n$ is an integer.

\section*{54}
(a): Note that
\[ \tan\left(\dfrac{\frac{\pi}{2}+L}{2}\right) \]
\[ =\dfrac{1-(\cos(\frac{\pi}{2}+L))}{\sin(\frac{\pi}{2}+L)} \]
\[ =\frac{1-(\cos\frac{\pi}{2}\cos L-\sin\frac{\pi}{2}\sin L)}{\sin\frac{\pi}{2}\cos L+\sin L\cos\frac{\pi}{2}} \]
\[ =\frac{1+\sin L}{\cos L} \]
Thus
\[ r\ln\left(\tan\left(\dfrac{\frac{\pi}{2}+L}{2}\right)\right) \]
\[ =r\ln\left(\frac{1+\sin L}{\cos L}\right) \]


(b): A reasonable domain is $-r\le x\le r$, as we are only considering points withint the radius $r$ of the globe.

\section*{56}
From the diagram, we can see that the area of the rectangle is $\sin\theta\cos\theta$. But
\[ \sin\theta\cos\theta=\frac{1}{2}\sin 2\theta \] Thus an equation modeling the area is
\[ y=\frac{1}{2}\sin 2\theta \]
which has an amplitude of $\frac{1}{2}$, a period of $\frac{2\pi}{2}=pi$, and thus maximum values of
\[ (\frac{\pi}{4}+\pi n,\frac{1}{2}) \] where $n$ is an integer. In the range of $0\le\theta\le 2\pi$, the value of $\theta$ which create the largest rectangle are
$\frac{\pi}{4}$ and $\frac{5\pi}{4}$.



\end{document}
